\begin{pastexam}[label={px:gk3-104-ds-02}]{高考三級資訊處理 104 年 第二題}
（申論；20 分）

樹（tree）是一個很常用的資料結構。一個樹是指一個沒有迴圈（cycle）的聯通圖（connected graph）。（每小題 10 分，共 20 分）
\begin{itemize}[leftmargin=2.2em]
\item[㈠] 證明：每個具有 $n$ 個節點（node）的樹，$n>1$，至少有 2 個分支度（degree）為 1 的節點。（分支度就是指有多少邊以此節點為端點。）
\item[㈡] 用前項結果證明：每個具有 $n$ 個節點的樹，$n>1$，恰好有 $n-1$ 個邊（edge）。
\end{itemize}
\end{pastexam}

\begin{pastsoln}
這裡的樹是圖論的樹（無根、無向），分支度就是圖的度數（資結書第~21.1~節）。 圖是簡單圖：沒有自迴圈、
兩點之間最多一條邊（有自迴圈或重邊就有迴圈，本來就不是樹）。 ㈡ 的結論就是資結書定理~21.3 的
(1)$\Rightarrow$(2)，這裡照題目要求用 ㈠ 來證。

\textbf{㈠}。 \textbf{證明。} 樹 $T$ 有 $n > 1$ 個節點又連通，所以至少有一條邊。 在 $T$ 的所有簡單路徑（節點
不重複的路徑）裡取邊數最多的一條 $P = \langle v_0, v_1, \ldots, v_k \rangle$；有邊，所以 $k \ge 1$，兩端
$v_0 \ne v_k$。 證明 $\deg(v_0) = 1$：$v_0$ 和 $v_1$ 相鄰，度數至少 $1$。 假設 $v_0$ 還有另一個鄰居
$u \ne v_1$：
\begin{itemize}
\item 若 $u$ 不在 $P$ 上，$\langle u, v_0, v_1, \ldots, v_k \rangle$ 是一條更長的簡單路徑，和 $P$ 最長矛盾。
\item 若 $u = v_j$（$2 \le j \le k$），則 $v_0, v_1, \ldots, v_j, v_0$ 是一個迴圈，和樹沒有迴圈矛盾。
\end{itemize}
所以 $\deg(v_0) = 1$。 同樣的論證用在另一端，$\deg(v_k) = 1$。 $v_0 \ne v_k$，所以至少有兩個分支度為 $1$
的節點。 \hfill$\blacksquare$

\textbf{㈡}。 \textbf{證明。} 對 $n$ 做數學歸納法（$n \ge 2$）。
\begin{itemize}
\item \textbf{基底 $n = 2$}：連通，所以兩點之間有邊；簡單圖兩點之間最多一條邊，所以恰好 $1 = n - 1$ 條。
\item \textbf{歸納}：設 $n \ge 3$，並假設每個 $n - 1$ 個節點的樹恰有 $n - 2$ 條邊。 由 ㈠，$n$ 個節點的樹 $T$
有一個分支度為 $1$ 的節點 $v$，令它唯一的邊是 $\{v, u\}$。 把 $v$ 和這條邊拿掉，得到 $T'$，有 $n - 1 \ge 2$
個節點。
  \begin{itemize}
  \item $T'$ 沒有迴圈：它是 $T$ 的子圖，$T$ 沒有迴圈。
  \item $T'$ 連通：任取 $T'$ 的兩點 $x, y$，它們在 $T$ 裡有一條簡單路徑。 這條路徑不會經過 $v$：路徑中間的節點
  要有兩條路徑上的邊（進來一條、出去一條），度數至少 $2$，而 $v$ 的度數是 $1$；$v$ 也不是端點（端點是
  $x, y \ne v$）。 所以這條路徑整條在 $T'$ 裡。
  \end{itemize}
  所以 $T'$ 是 $n - 1$ 個節點的樹，由歸納假設有 $n - 2$ 條邊。 $T$ 的邊比 $T'$ 多 $\{v, u\}$ 這一條，共
  $n - 1$ 條。
\end{itemize}
由歸納法，每個 $n > 1$ 個節點的樹恰好有 $n - 1$ 個邊。 \hfill$\blacksquare$

注意歸納步驟一定要說明拿掉葉節點以後\textbf{還是樹}（連通、無迴圈），才能套用歸納假設；這是本題的
得分重點。
\end{pastsoln}
